% ==========================================
% 第四章 随机变量的数字特征相关
% ==========================================

% --- 在本节开头局部定义新的引理颜色 ---
\makeatletter
% 定义 sec1_3 的专属颜色：橙色
\definecolor{LCCOrangeSec13}{RGB}{235, 120, 40} 
% 定义新的引理环境，完美继承 lcc@basestyle 和原有的计数器
\newtcolorbox[use counter from=引理]{引理Sec13}[1][]{
    lcc@basestyle,
    colframe=LCCOrangeSec13!70!white, colback=LCCOrangeSec13!2,
    interior style={top color=LCCOrangeSec13!8, bottom color=LCCOrangeSec13!1},
    frame style={left color=LCCOrangeSec13, right color=LCCOrangeSec13!60!white},
    title={引理~\thetcbcounter\if\relax\detokenize{#1}\relax\else\quad #1\fi}
}
\makeatother
% --------------------------------------

\vspace{1.5em}

\subsubsection*{1. 切比雪夫不等式 (Chebyshev's Inequality)}

\begin{引理Sec13}[切比雪夫不等式的意义]
	切比雪夫不等式给出了在随机变量分布未知的情况下，仅仅利用其期望和方差，就能估算出随机变量偏离其期望值的概率上限。
\end{引理Sec13}

设随机变量 $X$ 的数学期望为 $E(X) = \mu$，方差为 $D(X) = \sigma^2$，则对于任意给定的正数 $\varepsilon > 0$，有：
$$ P\{|X - \mu| \ge \varepsilon\} \le \frac{\sigma^2}{\varepsilon^2} $$
等价形式（对立事件）：
$$ P\{|X - \mu| < \varepsilon\} \ge 1 - \frac{\sigma^2}{\varepsilon^2} $$

\textbf{证明}\quad （以连续型随机变量为例，离散型同理可证）\\
设 $X$ 的概率密度函数为 $f(x)$，根据方差的定义：
$$ D(X) = \int_{-\infty}^{+\infty} (x - \mu)^2 f(x) dx $$
由于被积函数非负，我们可以缩小积分区间，去掉 $|x - \mu| < \varepsilon$ 的部分，不等号依然成立：
$$ D(X) \ge \int_{|x - \mu| \ge \varepsilon} (x - \mu)^2 f(x) dx $$
在积分区间 $|x - \mu| \ge \varepsilon$ 上，显然有 $(x - \mu)^2 \ge \varepsilon^2$，因此：
$$ 
\begin{aligned}
D(X) &\ge \int_{|x - \mu| \ge \varepsilon} \varepsilon^2 f(x) dx \\
&= \varepsilon^2 \int_{|x - \mu| \ge \varepsilon} f(x) dx \\
&= \varepsilon^2 P\{|X - \mu| \ge \varepsilon\}
\end{aligned}
$$
两边同除以 $\varepsilon^2$，即证得：
$$ P\{|X - \mu| \ge \varepsilon\} \le \frac{D(X)}{\varepsilon^2} = \frac{\sigma^2}{\varepsilon^2} $$


\vspace{1em}
\subsubsection*{2. 协方差与相关系数}

\textbf{协方差 (Covariance)}的定义与计算公式：
$$ \text{Cov}(X, Y) = E{[X - E(X)][Y - E(Y)]} $$
\textbf{计算公式推导：}
$$ 
\begin{aligned}
\text{Cov}(X, Y) &= E[XY - X \cdot E(Y) - Y \cdot E(X) + E(X)E(Y)] \\
&= E(XY) - E(X)E(Y) - E(Y)E(X) + E(X)E(Y) \\
&= E(XY) - E(X)E(Y)
\end{aligned}
$$

\textbf{相关系数 (Correlation Coefficient)}的定义：
$$ \rho_{XY} = \frac{\text{Cov}(X, Y)}{\sqrt{D(X)} \cdot \sqrt{D(Y)}} $$

\textbf{方差的展开公式：}
$$ D(X+Y) = D(X) + D(Y) + 2\text{Cov}(X, Y) $$
\textbf{证明：}
$$ 
\begin{aligned}
D(X+Y) &= E\{[(X+Y) - E(X+Y)]^2\} \\
&= E\{[(X - E(X)) + (Y - E(Y))]^2\} \\
&= E\{[X - E(X)]^2 + [Y - E(Y)]^2 + 2[X - E(X)][Y - E(Y)]\} \\
&= D(X) + D(Y) + 2\text{Cov}(X, Y)
\end{aligned}
$$

\textbf{协方差的性质及证明：}
\begin{itemize}
    \item 提取常数：$\text{Cov}(aX, bY) = ab \cdot \text{Cov}(X, Y)$ \\
    \textit{证明：} $\text{Cov}(aX, bY) = E[(aX - E(aX))(bY - E(bY))] = E[ab(X - E(X))(Y - E(Y))] = ab\text{Cov}(X, Y)$
    \item 分配律：$\text{Cov}(X_1 + X_2, Y) = \text{Cov}(X_1, Y) + \text{Cov}(X_2, Y)$
\end{itemize}

\textbf{不相关与独立的等价关系：}
$$ \rho_{XY} = 0 \iff X, Y \text{ 不相关 } \iff \text{Cov}(X, Y) = 0 \iff E(XY) = E(X) \cdot E(Y) $$

\begin{引理Sec13}[正态分布的特殊性]
	对于一般的随机变量，独立必定不相关，但不相关不一定独立。然而，对于\textbf{二维正态分布}，这两个概念是完全等价的。
\end{引理Sec13}

当 $(X, Y)$ 服从二维正态分布时：
$$ X, Y \text{ 相互独立 } \iff X, Y \text{ 不相关 } (\text{即 } \rho = 0) $$
\textbf{证明简述}\quad 若 $X, Y$ 不相关，即 $\rho = 0$，代入二维正态分布的概率密度展开式中：
$$ 
\begin{aligned}
f(x_1, x_2) &= \frac{1}{2\pi \sigma_1 \sigma_2} \exp\left\{ -\frac{1}{2} \left[ \frac{(x_1-\mu_1)^2}{\sigma_1^2} + \frac{(x_2-\mu_2)^2}{\sigma_2^2} \right] \right\} \\
&= \left[ \frac{1}{\sqrt{2\pi}\sigma_1} \exp\left(-\frac{(x_1-\mu_1)^2}{2\sigma_1^2}\right) \right] \cdot \left[ \frac{1}{\sqrt{2\pi}\sigma_2} \exp\left(-\frac{(x_2-\mu_2)^2}{2\sigma_2^2}\right) \right] \\
&= f_{X_1}(x_1) \cdot f_{X_2}(x_2)
\end{aligned}
$$
联合概率密度等于边缘概率密度的乘积，故 $X, Y$ 相互独立。反之，若独立，协方差必为 0，即 $\rho = 0$。


\vspace{1em}
\subsubsection*{3. 矩与协方差矩阵}

\textbf{矩 (Moments) 的定义}
\begin{itemize}
    \item $k$ 阶原点矩：$E(X^k)$
    \item $k$ 阶中心矩：$E\{[X - E(X)]^k\}$
    \item $k+l$ 阶混合矩：$E(X^k Y^l)$
    \item $k+l$ 阶混合中心矩：$E\{[X - E(X)]^k \cdot [Y - E(Y)]^l\}$
\end{itemize}

\textbf{二维正态分布与协方差矩阵}
设 $(X, Y) \sim N(\mu_1, \mu_2, \sigma_1^2, \sigma_2^2, \rho)$，则其相关系数 $\rho_{XY} = \rho$
其协方差矩阵 $C$ （或记为 $\Sigma$）定义为：
$$ C = \begin{pmatrix} \sigma_1^2 & \rho \sigma_1 \sigma_2 \\ \rho \sigma_1 \sigma_2 & \sigma_2^2 \end{pmatrix} $$

二维正态分布的概率密度函数有两种表示方法：

\textbf{(1) 矩阵形式 (更具一般性，易于向高维推广)：}
记均值向量 $\boldsymbol{\mu} = (\mu_1, \mu_2)^T$，随机向量 $\boldsymbol{X} = (x_1, x_2)^T$，则：
$$ f(x_1, x_2) = \frac{1}{(2\pi) \cdot |C|^{\frac{1}{2}}} \exp\left\{ -\frac{1}{2} (\boldsymbol{x} - \boldsymbol{\mu})^T C^{-1} (\boldsymbol{x} - \boldsymbol{\mu}) \right\} $$

\textbf{(2) 展开形式 (常用于代数计算)：}
$$ 
\begin{aligned}
f(x_1, x_2) &= \frac{1}{2\pi \sigma_1 \sigma_2 \sqrt{1-\rho^2}} \cdot \\
&\exp\left\{ -\frac{1}{2(1-\rho^2)} \left[ \frac{(x_1-\mu_1)^2}{\sigma_1^2} - 2\rho \frac{(x_1-\mu_1)(x_2-\mu_2)}{\sigma_1 \sigma_2} + \frac{(x_2-\mu_2)^2}{\sigma_2^2} \right] \right\}
\end{aligned}
$$